\begin{abstract}
On-policy distillation (OPD) transfers reasoning behavior through teacher feedback on student-generated trajectories. Common sampled-token methods assign this feedback independently to each token, although a local reasoning decision may span several tokens. We introduce Block3 Mean, which averages teacher--student log-probability differences within three-token blocks and optimizes a clipped objective with a joint block importance ratio and valid-length weighting. This gives adjacent tokens a shared learning signal without requiring an additional reward model or process annotations. We analyze how the update couples neighboring credit, changes gradient scale and alters clipping, and evaluate mathematical reasoning across teacher--student pairs and saved checkpoints. In the historical Qwen3-1.7B comparison at Step200, Block3 improves Avg@8 over sampled-token OPD by 8.90 percentage points on MATH500 and 5.42 on AMC23. The official Qwen3-8B-to-1.7B Instruct pair also has positive Step200 Avg@8 differences on all four evaluated benchmarks. Trajectory diagnostics connect these results to entropy, teacher--student overlap and local credit redistribution, while distinguishing numerical stability from generation quality. Our study makes feedback granularity an explicit design dimension in reasoning distillation.

\end{abstract}

\begin{figure}[htbp]
\centering
\includegraphics[width=\linewidth]{study_overview.png}
\caption{\textbf{From token-level to block-level updates.} The student generates a response and the frozen teacher scores the same prefixes. Token OPD keeps each token's feedback separate; Block3 Mean averages feedback within three adjacent tokens and shares it across the block. The implemented method also uses a joint block ratio and valid-length-weighted block loss. The two rows are alternative training runs, not summed objectives. Teacher feedback and old-policy probabilities are detached; only the current student receives gradients. A full three-token block is illustrated; masked tails use their valid length $n_j$. Both methods sample complete trajectories before updating.}
\label{fig:overview}
\end{figure}

\section{Introduction}
Transferring reasoning behavior from a capable teacher is an important approach to improving language models. On-policy distillation (OPD) lets the student generate its own responses and receive teacher feedback at the prefixes it actually visits \citep{gkd,minillm}. This brings supervision to the states induced by the learner, addressing the mismatch that can arise when training only on expert demonstrations \citep{dagger}. Teacher probabilities additionally provide feedback throughout a response, rather than only at its final answer.

An important choice remains: how should this feedback be assigned to the generated text? In sampled-token OPD, each token receives its own teacher--student log-probability difference \citep{revisiting}. This gives dense supervision, but density alone does not determine the appropriate unit of credit assignment. We ask: \textbf{can a short block of adjacent tokens serve as a useful feedback unit for on-policy distillation?}

Our intuition is that parts of a local reasoning decision, such as a short mathematical expression or transition, may span several tokens. Their individual probability differences can reflect both the underlying choice and its particular textual realization. Sharing a local signal could reduce sensitivity to such fluctuations while retaining finer granularity than a response-level reward.

Motivated by this perspective, we introduce \method{} (Figure~\ref{fig:overview}). The student samples a response, and a frozen teacher scores the same prefixes. We partition the response into disjoint three-token blocks, average the distillation advantages within each block, and use the shared advantage with a joint block importance ratio and a valid-length-weighted objective. This reuses the existing teacher scores without adding a reward model or process annotations. Training still collects complete trajectories before each update.

We evaluate on MATH500, AIME24, AIME25 and AMC23, reporting each benchmark separately. Within each matched comparison, the two methods share initialization, data, prompting and evaluation settings. The historical Qwen3-1.7B comparison improves Step200 Avg@8 by 8.90 points on MATH500 and 5.42 on AMC23; the official Qwen3-8B-to-1.7B Instruct pair has positive Step200 differences on all four benchmarks.

Our contributions are threefold. First, we formulate a block-level OPD update that combines local advantage sharing with a joint importance ratio. Second, we derive its cross-token credit terms and gradient scale, and give conditions under which averaging improves a local signal estimate. Third, we provide checkpoint-level comparisons and trajectory diagnostics covering entropy, token overlap, credit redistribution and generation behavior. Together, these analyses relate the method's observed behavior to its optimization structure.

\section{Related Work}
\label{sec:related}
GKD trains on student-generated sequences and permits multiple choices of divergence; on-policy sampling therefore does not uniquely require reverse KL \citep{gkd}. MiniLLM studies reverse-KL-based distillation for language generation \citep{minillm}. Our baseline follows the sampled-token OPD implementation in Revisiting OPD, which studies weaknesses of local distillation signals and practical corrections \citep{revisiting}. We do not equate this baseline with an unbiased estimator of the full sequence-level KL gradient.

Rethinking OPD examines teacher--student compatibility and training dynamics, motivating our entropy, top-token overlap and position diagnostics \citep{rethinking}. Blockwise policy-drift gating aggregates old--current student log-probability shifts into detached, mean-normalized weights while retaining the token objective \citep{bpdg}. FiRe-OPD filters trajectories and reweights tokens; TIP selects informative token supervision \citep{fireopd,tip}. Our intervention instead shares teacher feedback and changes the PPO ratio itself. R2-OPD filters spans using independent reasoning-progress estimates \citep{r2opd}; GSPO uses a length-normalized sequence ratio, unlike our product of token ratios \citep{gspo}. These precedents preclude a generic claim of first introducing blockwise, selective, or multi-step OPD. We have not benchmarked their best recipes head to head.

\section{Local Credit Assignment}
\label{sec:method}
\subsection{From Individual Feedback to Local Credit}
Token-wise updates follow each teacher preference independently; block-wise updates aggregate the preferences and update the expression's components with one signal. In Figure~\ref{fig:credit}, averaging $(2,-1,2)$ gives positive feedback at all three positions. This couples adjacent updates. A fixed-width block is an operational unit, not an annotated reasoning step.

The student generates a complete response; the frozen teacher scores its prefixes; detached feedback is aggregated locally; and gradients update only the current student. No new teacher responses, reward model or process labels are needed. Teacher scoring can process the generated sequence in a batched forward pass, but grouping does not reduce the number of conditional probabilities scored.

\subsection{Block-Level On-Policy Distillation}
\paragraph{Shared feedback and joint ratio.}
Group consecutive response tensor positions into disjoint blocks $B_j$ of width $k$, without packing across mask holes. Below, sums within a block select valid positions, $n_j$ is their count, and $T=\sum_j n_j$. Empty blocks have zero weight. The ideal real-arithmetic reduction, omitting the implementation's $10^{-8}$ denominator stabilizer, is
\begin{align}
 \bar a_j&=\frac{1}{\max(1,n_j)}\sum_{t\in B_j}a_t,
 &R_j(\theta)&=\exp\left(\sum_{t\in B_j}\log r_t(\theta)\right),\label{eq:block}\\
 \mathcal L_{\mathrm{block}}&=-\frac1T\sum_{j:n_j>0}n_j\mathcal C(R_j,\sg(\bar a_j)),
 &\mathcal L_{\mathrm{token}}&=-\frac1T\sum_{t=1}^T\mathcal C(r_t,\sg(a_t)).\label{eq:loss}
\end{align}
\method{} uses $k=3$. Its joint ratio $R_j=\prod_{t\in B_j}r_t$ measures the product of policy changes across the sampled block, rather than their average. Valid-length weighting gives a two-token tail weight $2/3$ relative to a full block. The code implements the same reduction using fractional block masks $n_j/k$; Appendix~\ref{app:proofs} gives the exact stabilized denominator and micro-batch convention.

\paragraph{Optimization.}
Each collected batch supplies shared advantages and current-to-old block ratios for the clipped update. This uses sampled blocks, not enumeration of a $|V|^k$ distribution or optimizer steps after every three tokens. Relative to Token OPD, the intervention changes advantage sharing, ratio granularity and loss reduction together.

\begin{figure}[t]
\centering
\includegraphics[width=.94\linewidth]{credit_assignment_concept.png}
\caption{Illustrative feedback sharing, not an observed rollout. Averaging $(2,-1,2)$ gives $(1,1,1)$, reversing the middle token's local signal. The implemented method additionally changes ratios and loss reduction. A local sign reversal does not by itself establish an error in task-level credit assignment.}
\label{fig:credit}
\end{figure}

\subsection{Understanding the Block Update}
\paragraph{Cross-token credit.}
Consider equal-sized blocks, $T=kJ$, at the old policy, where every ratio is one and clipping is inactive. Let $g_t=\nabla_\theta\log\pi_\theta(y_t\mid h_t)$. The policy-gradient components satisfy
\begin{align}
 \nabla\mathcal L_{\mathrm{token}}&=-\frac1T\sum_t a_tg_t,\nonumber\\
 \nabla\mathcal L_{\mathrm{block}}&=-\frac1T\sum_j\sum_{i\in B_j}\sum_{t\in B_j}a_i g_t
 =\nabla\mathcal L_{\mathrm{token}}-\frac1T\sum_j\sum_{\substack{i,t\in B_j\\i\ne t}}a_i g_t.
 \label{eq:cross}
\end{align}
The additional terms send token $i$'s signal to token $t$'s score gradient. They can align or conflict with the baseline; their sign alone does not tell us whether downstream mathematical reasoning benefits.

\begin{proposition}[Equal feedback does not imply equal updates]
At $\theta=\theta_{\mathrm{old}}$, with inactive clipping, full blocks of size $k$, and $a_t=a$ for all valid tokens, both policy losses equal $-a$, but
$\nabla\mathcal L_{\mathrm{block}}=k\nabla\mathcal L_{\mathrm{token}}$.
\end{proposition}
Substituting $\bar a_j=a$ in Equation~\ref{eq:loss} proves the statement. The factor concerns the policy-loss derivative, not the subsequent Adam update after gradient clipping and other terms. It reveals a scale change even when all local feedback agrees.

\paragraph{Clipping geometry.}
Away from the old policy, the joint ratio also changes which updates are clipped. Three token ratios of 1.1 produce a block ratio of $1.1^3=1.331$: each token can remain inside a conventional 0.2 clipping band while the block leaves it. Learning-rate rescaling alone therefore cannot generally recover the token-level update.

\paragraph{When averaging could help.}
To separate the averaging hypothesis from the implementation above, consider a hypothetical token-ratio, token-normalized intervention using $Wa$ instead of $a$, where $W$ is deterministic and each row sums to one. Suppose, only as a local signal model, $a=\mu+\varepsilon$, $\E\varepsilon=0$, and $\Cov(\varepsilon)=\Sigma$. Then
\begin{equation}
 \E\|Wa-\mu\|_2^2=\|(W-I)\mu\|_2^2+\operatorname{tr}(W\Sigma W^\top).
 \label{eq:mse}
\end{equation}
Averaging helps this estimator when noise reduction exceeds the bias introduced by combining different latent signals. With a constant signal and independent equal-variance noise within a block, the per-position variance becomes $\sigma^2/k$; correlated noise or heterogeneous signals can remove this benefit. This explains the motivation for short-range sharing, but is a statement about signal estimation, not a guarantee of higher task reward. The implemented objective additionally changes ratios and reduction. Appendix~\ref{app:proofs} provides the derivation, a counterexample to unconditional smoothing benefit, and a restricted interpretation in which block cross terms recover within-block future credit.


\section{Experimental Protocol}
\label{sec:protocol}
\paragraph{Evidence units.}
We treat each model pair, prompt version, data pool, grader, and training seed as a separate experimental stratum. Within a stratum, Token and Block3 start independently from the same original student. They do not continue from each other's weights. For each evaluated checkpoint, we report the four mathematical benchmarks separately. Results from older small evaluations and different grading implementations are exploratory evidence, not extra repetitions of the main comparison. Appendix~\ref{app:history} inventories all located experimental series, including interrupted attempts and missing evidence.

\paragraph{Matched 4B Base training.}
The student is Qwen3-4B-Base, not the similarly named instruction-tuned model \citep{qwen3}. The teacher is the public Qwen3-4B-Base-GRPO checkpoint identified in Appendix~\ref{app:config}. Both arms retain the same DAPO-derived pool and frozen implementation. The local parquet repeats a 17,917-row question--answer sequence 100 times, totaling 1,791,700 rows with 17,237 exact unique question strings. Each of 200 updates uses four prompts and eight responses per prompt: 800 prompt occurrences, 783 exact distinct questions and 6,400 trajectories per arm. We use training seed 21, learning rate $2\times10^{-6}$, one PPO epoch per update, and four A100 GPUs. Complete states are retained at Steps 50, 100, 150, and 200. One normalized-string MATH500 overlap is encountered once per arm; Appendix~\ref{app:config} discloses pool/exposure audits and a separate sensitivity analysis.

Both arms use temperature 1.0, top-$p=0.9$, a 2,048-token prompt limit and a 16,384-token response cap. The generation seed is derived per request from the global seed, step, source identity, sample index and turn; it is not reset to the same seed for all eight responses. A post-training audit verifies all 6,400 paired input identities and sampling settings. Equal input streams do not imply equal later responses: diverging policies are supposed to generate different on-policy data. Training, generation and teacher-log-probability micro-batches are specified separately in Appendix~\ref{app:config}.

\paragraph{Prompt control.}
The current Base-model comparison bypasses ChatML. Training and evaluation use exactly the same function:
\begin{quote}\small
\texttt{[question]}\\
\texttt{Please solve the problem step by step and put the final answer in}\ \verb|\boxed{}.|\\
\texttt{Solution:}
\end{quote}
Blank lines separate these components, with a final newline after \texttt{Solution:}. Thinking is explicitly disabled; ordinary step-by-step reasoning remains permitted. Input token IDs, rather than only prompt text, are checked after checkpoint merging. The EOS token is 151643, with no additional stop-token IDs. Older Qwen and Llama experiments deliberately retain their historical prompting and are labeled accordingly, rather than silently reclassified as this matched completion protocol.

\paragraph{Evaluation and endpoints.}
The additional official Instruct pair retains the same data, optimizer settings, seed21, 200 updates and four checkpoints. Its prompt has the same boxed-answer instruction but uses native Qwen chat, explicitly disables thinking, and omits \texttt{Solution:}. EOS is 151645 with stop IDs $\{151645,151643\}$. The original student is also evaluated. Appendix~\ref{app:instruct} specifies its model revisions, input checks and recovery history; it is not a repeat of the older 1.7B/public-GRPO stratum.

Each checkpoint is evaluated on MATH500 (500 questions), AIME24 (30), AIME25 (30), and AMC23 (83), for 643 questions and 5,144 responses. Evaluation uses eight sampling seeds, 21 through 28, temperature 1.0, top-$p=0.9$, and the same 16,384-token cap for both arms. Four independent tensor-parallel-size-one workers run the inference workload. We retain the pinned historical mathematical grader and every raw response, input token sequence, sampling seed and engine stopping reason. The inference memory fraction is 0.9, distinct from the 0.6 training rollout setting. CPU checkpoint merging does not modify the training weights.

For question $i$, let $c_i\in\{0,\ldots,8\}$ count correctly graded responses. We compute
\begin{equation}
 \avg=\frac{1}{N}\sum_{i=1}^N\frac{c_i}{8},\qquad
 \pass=\frac{1}{N}\sum_{i=1}^N\mathbf{1}[c_i>0].
 \label{eq:metrics}
\end{equation}
The first measures mean sampled accuracy; the second measures success with eight attempts. Neither is a vote over eight answers. We use the scheduled final Step200 as the primary reporting endpoint, a retrospective reporting choice rather than a preregistration. Intermediate checkpoints characterize the trajectory; none is selected as a method's best result. We do not aggregate the four benchmarks. Where reported, intervals resample paired questions within a benchmark, conditional on the trained checkpoints. They do not measure variation across independent training seeds.

\paragraph{Behavioral and numerical diagnostics.}
We retain student and teacher entropy, signed and absolute entropy differences, top-16 overlap and overlap probability mass, overlap-token advantage, pre-clipping gradient norm, policy-gradient loss, response length and cap incidence. Position heatmaps summarize on-policy prefixes visited during training, not a fixed evaluation set. Block-specific sign disagreement and advantage redistribution statistics supplement these metrics. Missing literal \verb|\boxed| markers and engine length stops are reported separately: neither is synonymous with mathematical error, and the absence of a marker is only a specific format-failure indicator. Definitions and mask handling are given in Appendix~\ref{app:diagnostics}.

\section{Results}
\label{sec:results}
\paragraph{The 4B Base endpoint does not establish a general advantage.}
At Step200, Block3's Avg@8 differences from Token are $-0.52$, $-2.92$, $+0.42$ and $-3.01$ percentage points on MATH500, AIME24, AIME25 and AMC23, respectively. These are benchmark-specific descriptive differences, not four independent training replications. All pointwise 95\% paired-question intervals cross zero; the experiment establishes neither a statistically resolved advantage nor equivalence. Pass@8 tells a complementary story: AIME24 has lower average success but one more question solved at least once by Block3 (33.33\% versus 30.00\%). The two metrics must not be collapsed into a single claim of improvement.
% Generated from accepted results; requires booktabs.
\begin{table*}[t]
\centering
\scriptsize
\setlength{\tabcolsep}{4pt}
\caption{Qwen3-4B paired Step200 evaluation (seed21; queue COMPLETE).}
\label{tab:qwen4endpoint}
\begin{tabular}{lrrrrrr}
\toprule
 & \multicolumn{3}{c}{Avg@8} & \multicolumn{3}{c}{Pass@8} \\
\cmidrule(lr){2-4}\cmidrule(lr){5-7}
Benchmark & Block3 & Token & $\Delta$ [95\% CI] & Block3 & Token & $\Delta$ [95\% CI] \\
\midrule
MATH500 & 75.40 & 75.92 & -0.52 [-1.57, 0.55] & 90.60 & 91.00 & -0.40 [-2.40, 1.60] \\
AIME24 & 13.33 & 16.25 & -2.92 [-7.08, 0.83] & 33.33 & 30.00 & +3.33 [-6.67, 13.33] \\
AIME25 & 13.75 & 13.33 & +0.42 [-2.92, 3.75] & 30.00 & 30.00 & +0.00 [-13.33, 13.33] \\
AMC23 & 46.23 & 49.25 & -3.01 [-6.17, 0.30] & 74.70 & 77.11 & -2.41 [-10.84, 6.02] \\
\bottomrule
\end{tabular}
\par\smallskip\footnotesize Scores are percentages; $\Delta$ is Block3 minus Token in percentage points. Brackets: pointwise 95\% paired-question bootstrap within each benchmark, conditional on one training pair, not a multi-training-seed CI. NA: unavailable.
\end{table*}

\IfFileExists{\qwenfigures/qwen4_accuracy.pdf}{%
\begin{figure}[t]
\centering
\includegraphics[width=\linewidth]{\qwenfigures/qwen4_accuracy.pdf}
\caption{Current matched 4B comparison at all four saved checkpoints. Avg@8 is reported separately for each benchmark, with task-specific vertical ranges shared by both methods. Lines connect measured checkpoints, not additional evaluations. One training seed per method.}
\label{fig:qwen4accuracy}
\end{figure}
}{}

\paragraph{Intermediate gains reverse across checkpoints.}
The endpoint does not describe the entire trajectory. At Step150, Block3 improves MATH500 Avg@8 by 1.85 points (pointwise interval $[0.70,3.00]$) and AIME24 by 2.92 points ($[0.42,5.83]$), but reduces AIME25 by 4.17 points. At Step100, the MATH500 difference instead equals $-1.40$ points ($[-2.57,-0.25]$); at Step50, only AMC23 has a positive Avg@8 difference. These conditional, pointwise intervals do not correct for checkpoint/task selection or replicate training seeds. The results show checkpoint-dependent differences, not a consistent improvement or a validated best-checkpoint selection rule.

\paragraph{A new official Instruct pair has small positive mean-accuracy differences.}
For Qwen3-8B teaching Qwen3-1.7B, Block3's Step200 Avg@8 differences are $+0.10$, $+2.50$, $+0.83$ and $+0.30$ points (Table~\ref{tab:qwen17endpoint}). All pointwise 95\% paired-question intervals include zero. Pass@8 differences are $+0.20$, $-6.67$, $0.00$ and $-1.20$ points: higher average accuracy is not uniformly better eight-attempt coverage. Both methods exceed the original student's four Avg@8 scores, but both reduce its MATH500 Pass@8 from 91.20\%. Checkpoints also reverse: Step50 favors Token on three Avg@8 tasks, and Step150 favors Token on AIME24. Appendix~\ref{app:instruct} retains all nine model views and behavior metrics. This is a separate matched stratum, not a new seed of the older 1.7B experiment or proof that model scale or prompting caused the difference.
\begin{table}[t]
\centering\small
\caption{Official Qwen3-8B to Qwen3-1.7B Instruct: Step200 scores (\%) and Block3 minus Token differences (percentage points).}
\label{tab:qwen17endpoint}
\begin{tabular}{lrrrr}
\toprule
Benchmark & Original & Token & Block3 & $\Delta$ [95\% CI] \\
\midrule
\multicolumn{5}{c}{Avg@8} \\
MATH500 & 71.45 & 72.12 & 72.22 & +0.10 [-1.18, 1.38] \\
AIME24 & 10.00 & 14.17 & 16.67 & +2.50 [-2.50, 8.33] \\
AIME25 & 10.42 & 12.08 & 12.92 & +0.83 [-3.75, 5.83] \\
AMC23 & 40.51 & 43.07 & 43.37 & +0.30 [-3.16, 3.92] \\
\midrule
\multicolumn{5}{c}{Pass@8} \\
MATH500 & 91.20 & 89.60 & 89.80 & +0.20 [-2.00, 2.40] \\
AIME24 & 26.67 & 43.33 & 36.67 & -6.67 [-20.00, 6.67] \\
AIME25 & 23.33 & 30.00 & 30.00 & +0.00 [-16.67, 16.67] \\
AMC23 & 69.88 & 75.90 & 74.70 & -1.20 [-9.64, 7.23] \\
\bottomrule
\end{tabular}
\par\smallskip\footnotesize Original denotes the untrained-for-this-experiment Instruct student, not a Base model. Brackets are pointwise 95\% paired-question bootstrap intervals (10,000 resamples); one training seed, no correction for checkpoint or task selection.
\end{table}


\paragraph{Earlier positive results coexist with failed transfers.}
Table~\ref{tab:historyendpoints} reports the historical Step200 endpoints within their original protocol strata. The 1.7B comparison shows sizeable differences in favor of Block3 on all four Avg@8 benchmarks. The 1.5B DeepSeek/JustRL setting gives much smaller differences on MATH500 and AMC23 and larger ones on AIME, rather than a uniformly negative transfer. By contrast, the non-thinking 0.6B Block3 model nearly loses all measured mathematical accuracy while its Token control improves over the initial student. The Llama pair also favors Token. Model scale alone cannot explain this pattern: teacher adaptation, initial policy, prompt protocol, seed handling and implementation versions differ between strata. These experiments support a statement of conditional behavior, not a causal scaling law.

\paragraph{Longer blocks and numerical health.}
The historical Block10 runs exhibit severe generation degradation and recorded Step200 failures despite finite optimizer diagnostics. Their maximum pre-clipping gradient norms are nevertheless large (867.50 and 865.21); ``finite'' must not be read as ``well behaved.'' The older built-in grader records zeros on all four tasks, but its AMC23 compatibility caveat prevents treating that particular zero as independent proof of no capability (Appendix~\ref{app:history}). The remaining evidence still rules out using the mere absence of NaN or Inf gradients as a sufficient acceptance criterion. It does not isolate which of joint-ratio clipping, credit redistribution, prompting or rollout drift caused the failure. Randomly shifted blocks and sliding windows were also evaluated at several checkpoints. We report every retained result rather than identifying a winning window from an adaptive search without a new matched control.

\paragraph{Prompt defects can precede an algorithmic intervention.}
The first 4B attempt exhibited severe repetition and frequent cap hits before any optimizer update. That behavior cannot be attributed to Block3 training. The repaired comparison uses a completion prompt without ChatML and independent per-request seeds; because these fixes were introduced together, the resulting comparison is not a single-factor estimate of the effect of removing chat delimiters. Both methods nevertheless receive the same repaired protocol. Their training-cap incidences are 263/6,400 (4.11\%) for Token and 247/6,400 (3.86\%) for Block3. A small difference in this training statistic is not evidence of a general improvement in answer formatting or benchmark accuracy. Evaluation-specific format and length-stop counts are retained separately.

\paragraph{Credit disagreement without a Block10-like collapse.}
Figure~\ref{fig:qwen4training} shows low late-training entropy and largely reduced gradient norms in the current 4B pair. The median Block3 sign-disagreement rate over the 200 update summaries is 15.57\%, or 6.68\% when weighted by absolute local advantage. Its median relative redistribution magnitude is 1.148; this statistic is not a probability and can exceed one. Nonzero redistribution therefore coexists with the absence of the severe historical Block10 behavior, and cannot by itself diagnose collapse. Token's size-one redistribution statistics are zero by definition, not independent evidence of better credit assignment. Appendix~\ref{app:positionfigures} reports every-step position maps and their coverage limits.


\section{Discussion and Limitations}
\label{sec:limits}
Local aggregation presents a bias--variance trade-off, but the historical implementation presents more than one such change at once. A block shares rewards, multiplies importance ratios, and changes the reduction denominator. The gradient calculation in Section~\ref{sec:method} rules out attributing every gain solely to semantic agreement between adjacent tokens. A decisive follow-up would hold the token ratio and token normalization fixed while changing only the averaging operator, then separately vary ratio granularity and learning-rate scale. We have not run that complete factorial study and do not report it as established evidence.

The evidence is also limited by training replication. A single trained seed with thousands of evaluation responses is still one optimization trajectory; eight evaluation samples do not create eight independent training runs. AIME results have especially coarse question-level resolution: one additional solved question changes Pass@8 by 3.33 percentage points. Window variants, checkpoints and pilot weights were explored adaptively. Their best observed scores cannot be interpreted as unbiased estimates of a selected method's advantage.

Prompt diagnostics reveal another limitation of transferring an implementation across models. Base and instruction-tuned models need not react similarly to role delimiters or thinking instructions. A difference present before the first optimizer update cannot have been caused by the block loss. Conversely, a short prompt probe does not establish that the new protocol eliminates repetition throughout training. We therefore keep prompt validation, optimizer health, generation quality and benchmark correctness as distinct checks. Position-wise entropy is useful for locating changes, but on-policy occupancy changes across training, so such heatmaps alone do not identify a causal mechanism.

Finally, these are comparisons against a sampled-token baseline, not every recent OPD method. We establish neither efficiency gains, reduced teacher computation nor universal format correctness. The 4B completion table supports method-to-method differences, without a matched original-student result. The new Instruct pair includes that control, but its shared data pool is not an independent dataset replication. More training seeds, component ablations and new task distributions remain necessary.

\section{Conclusion}
We introduced Block3 Mean, an on-policy distillation update that combines shared feedback over adjacent tokens with a joint block ratio and valid-length weighting. The method reuses teacher scores while changing the unit through which they influence the student. Experiments show substantial gains in the historical Qwen3-1.7B setting and positive final Avg@8 differences for the official 8B-to-1.7B Instruct pair, alongside checkpoint reversals and settings that favor token-level updates. The analysis explains how the objective introduces cross-token credit, scale changes and a different clipping geometry. Together with the retained trajectory diagnostics, these results identify local credit assignment as a concrete design choice for reasoning distillation and provide a basis for separating its optimization components.

\FloatBarrier
\label{mainend}

\section*{Reproducibility Statement}
The loss definitions and derivation appear in Section~\ref{sec:method} and Appendix~\ref{app:proofs}. Appendices~\ref{app:config} and~\ref{app:instruct} record the paired protocols; Appendix~\ref{app:history} distinguishes historical versions and evidence gaps. Both matched pairs retain complete states, all trajectories and input checks. Tables and figures use provenance-linked summaries. This does not imply equivalent retention for every pilot or an already public anonymous artifact release.

\section*{Ethics Statement}
The study uses mathematical reasoning benchmarks and pretrained language models rather than human-subject data collection. Better imitation of a teacher does not by itself establish truthfulness or safety. We retain negative outcomes and avoid converting formatting success, finite gradients or benchmark-specific gains into general capability claims. Any distributed artifacts must respect the upstream model and dataset licenses and omit credentials, identifying infrastructure paths and private checkpoints.

\section*{AI Use Statement}
Generative AI assistants were used to discuss and refine hypotheses; assist with experimental design, mathematical arguments and implementation; debug and monitor experiments; search and summarize related work; process evidence and interpret results; and draft, revise and translate this manuscript and its figures. The student and teacher models additionally generate the on-policy training trajectories that constitute the method itself. Automated tests and artifact audits check specified implementation and data contracts; they do not substitute for independent replication or the authors' responsibility for the final submission. No claim is made that the research, proofs, writing or translation were produced without AI assistance.
